%% Introduction — body only (no \section{} header)
%% Compiled via main.tex %% Introduction — body only (no \section{} header)
%% Compiled via main.tex %% Introduction — body only (no \section{} header)
%% Compiled via main.tex %% Introduction — body only (no \section{} header)
%% Compiled via main.tex \input{sections/introduction}

When governments closed schools to contain COVID-19, they traded an immediate health risk for an educational one. More than 1.6 billion learners across over 190 countries experienced school interruptions during the pandemic \citep{unesco2021year}. The central empirical question is how much those interruptions cost students in learning. The answer is large enough to matter for a generation of workers and for the public investments required to recover.

The evidence from systematic reviews is clear: pandemic-related closures reduced student achievement by economically meaningful amounts, and the losses fell disproportionately on disadvantaged children. \citet{betthauser2023systematic} synthesize 42 studies across 15 countries and document a mean learning deficit of 0.14 standard deviations (SD)---equivalent to roughly one-third of a typical school year---with consistently larger losses among socio-economically disadvantaged students. \citet{patrinos2022analysis} review 36 rigorous evaluations and report a mean shortfall of 0.17 SD, approximately half a grade level, alongside systematic widening of socio-economic gaps. \citet{delacruz2024learning} extend the evidence to 56 studies drawn primarily from developing countries and estimate that each full year of school closure cost approximately 1.1 years of learning; they also identify in-person and phone-based tutoring as effective remediation tools. Together these reviews establish that COVID-19 generated large, unequal, and only partially recovered losses worldwide.

Country-level evidence has accumulated rapidly in middle-income settings. Studies in Mexico \citep{hevia2022estimation, alasino2024learning}, South Africa \citep{ardington2021covid}, India \citep{guariso2023impact, singh2024covid}, and Brazil \citep{lichand2022impacts} use in-person assessments of large student samples to document losses of comparable magnitude, while revealing important variation in severity tied to the duration of closure and the quality of remote instruction. Ukraine, an upper-middle-income country, is absent from this literature despite operating a nationally representative primary-school monitoring programme with assessments administered under standardized conditions.

I use two rounds of Ukraine's State Monitoring of Quality of Primary Education
\citep{uceqa2025} to estimate pandemic-related learning losses in mathematics and reading. The programme assessed fourth-grade students in May 2018---before any school disruption---and again in May 2021, after two academic years marked by extended closures and shifts between remote and hybrid instruction. I compare cross-cohort achievement distributions across waves after controlling for region, school type, and student socioeconomic background.

I find that average mathematics scores fell by 0.09 SD between the two cohorts, a decline equivalent to roughly one-fifth of a typical school year \citep{hill2008empirical}. Reading scores fell by a similar magnitude, 0.09--0.10 SD. These losses are smaller than the cross-country averages reported in \citet{betthauser2023systematic} and \citet{patrinos2022analysis}, and I document the mechanisms behind this gap: nearly half of students experienced only one to two months of distance learning, while 13 percent were exposed for three months or longer. I use teacher questionnaires linked to student records to characterize heterogeneity in instructional practices during remote periods.

This paper makes three contributions. First, it provides the first nationally representative estimates of COVID-19 learning losses in Ukrainian primary education. The State Monitoring programme administered identical test instruments under standardized field conditions in both waves and anchored scores to a common scale, which allows direct cross-cohort comparison without the equating assumptions that complicate many international comparisons. Ukraine's experience---an upper-middle-income country with moderately long but variable closure durations---extends the middle-income evidence base to a regional setting that has received little systematic attention. Second, because the design compares cross-sections rather than following the same students, I document sample comparability across waves and assess sensitivity to differential test participation and cohort compositional change. Third, the monitoring programme links student scores to classroom-level teacher questionnaires; I use these data to document heterogeneity in distance-learning exposure and instructional practices, connecting variation in remote-instruction quality to variation in measured losses.

The paper proceeds as follows. Section~\ref{sec:background} describes Ukraine's school closure policies and the remote-learning environment. Section~\ref{sec:data} introduces the monitoring programme, the sample, and the main variables. Section~\ref{sec:results} presents the cross-cohort achievement comparisons. Section~\ref{sec:robustness} reports sensitivity analyses addressing participation and compositional change. Section~\ref{sec:heterogeneity} examines heterogeneity by distance-learning exposure and socioeconomic background. Section~\ref{sec:conclusion} concludes with policy implications.


When governments closed schools to contain COVID-19, they traded an immediate health risk for an educational one. More than 1.6 billion learners across over 190 countries experienced school interruptions during the pandemic \citep{unesco2021year}. The central empirical question is how much those interruptions cost students in learning. The answer is large enough to matter for a generation of workers and for the public investments required to recover.

The evidence from systematic reviews is clear: pandemic-related closures reduced student achievement by economically meaningful amounts, and the losses fell disproportionately on disadvantaged children. \citet{betthauser2023systematic} synthesize 42 studies across 15 countries and document a mean learning deficit of 0.14 standard deviations (SD)---equivalent to roughly one-third of a typical school year---with consistently larger losses among socio-economically disadvantaged students. \citet{patrinos2022analysis} review 36 rigorous evaluations and report a mean shortfall of 0.17 SD, approximately half a grade level, alongside systematic widening of socio-economic gaps. \citet{delacruz2024learning} extend the evidence to 56 studies drawn primarily from developing countries and estimate that each full year of school closure cost approximately 1.1 years of learning; they also identify in-person and phone-based tutoring as effective remediation tools. Together these reviews establish that COVID-19 generated large, unequal, and only partially recovered losses worldwide.

Country-level evidence has accumulated rapidly in middle-income settings. Studies in Mexico \citep{hevia2022estimation, alasino2024learning}, South Africa \citep{ardington2021covid}, India \citep{guariso2023impact, singh2024covid}, and Brazil \citep{lichand2022impacts} use in-person assessments of large student samples to document losses of comparable magnitude, while revealing important variation in severity tied to the duration of closure and the quality of remote instruction. Ukraine, an upper-middle-income country, is absent from this literature despite operating a nationally representative primary-school monitoring programme with assessments administered under standardized conditions.

I use two rounds of Ukraine's State Monitoring of Quality of Primary Education
\citep{uceqa2025} to estimate pandemic-related learning losses in mathematics and reading. The programme assessed fourth-grade students in May 2018---before any school disruption---and again in May 2021, after two academic years marked by extended closures and shifts between remote and hybrid instruction. I compare cross-cohort achievement distributions across waves after controlling for region, school type, and student socioeconomic background.

I find that average mathematics scores fell by 0.09 SD between the two cohorts, a decline equivalent to roughly one-fifth of a typical school year \citep{hill2008empirical}. Reading scores fell by a similar magnitude, 0.09--0.10 SD. These losses are smaller than the cross-country averages reported in \citet{betthauser2023systematic} and \citet{patrinos2022analysis}, and I document the mechanisms behind this gap: nearly half of students experienced only one to two months of distance learning, while 13 percent were exposed for three months or longer. I use teacher questionnaires linked to student records to characterize heterogeneity in instructional practices during remote periods.

This paper makes three contributions. First, it provides the first nationally representative estimates of COVID-19 learning losses in Ukrainian primary education. The State Monitoring programme administered identical test instruments under standardized field conditions in both waves and anchored scores to a common scale, which allows direct cross-cohort comparison without the equating assumptions that complicate many international comparisons. Ukraine's experience---an upper-middle-income country with moderately long but variable closure durations---extends the middle-income evidence base to a regional setting that has received little systematic attention. Second, because the design compares cross-sections rather than following the same students, I document sample comparability across waves and assess sensitivity to differential test participation and cohort compositional change. Third, the monitoring programme links student scores to classroom-level teacher questionnaires; I use these data to document heterogeneity in distance-learning exposure and instructional practices, connecting variation in remote-instruction quality to variation in measured losses.

The paper proceeds as follows. Section~\ref{sec:background} describes Ukraine's school closure policies and the remote-learning environment. Section~\ref{sec:data} introduces the monitoring programme, the sample, and the main variables. Section~\ref{sec:results} presents the cross-cohort achievement comparisons. Section~\ref{sec:robustness} reports sensitivity analyses addressing participation and compositional change. Section~\ref{sec:heterogeneity} examines heterogeneity by distance-learning exposure and socioeconomic background. Section~\ref{sec:conclusion} concludes with policy implications.


When governments closed schools to contain COVID-19, they traded an immediate health risk for an educational one. More than 1.6 billion learners across over 190 countries experienced school interruptions during the pandemic \citep{unesco2021year}. The central empirical question is how much those interruptions cost students in learning. The answer is large enough to matter for a generation of workers and for the public investments required to recover.

The evidence from systematic reviews is clear: pandemic-related closures reduced student achievement by economically meaningful amounts, and the losses fell disproportionately on disadvantaged children. \citet{betthauser2023systematic} synthesize 42 studies across 15 countries and document a mean learning deficit of 0.14 standard deviations (SD)---equivalent to roughly one-third of a typical school year---with consistently larger losses among socio-economically disadvantaged students. \citet{patrinos2022analysis} review 36 rigorous evaluations and report a mean shortfall of 0.17 SD, approximately half a grade level, alongside systematic widening of socio-economic gaps. \citet{delacruz2024learning} extend the evidence to 56 studies drawn primarily from developing countries and estimate that each full year of school closure cost approximately 1.1 years of learning; they also identify in-person and phone-based tutoring as effective remediation tools. Together these reviews establish that COVID-19 generated large, unequal, and only partially recovered losses worldwide.

Country-level evidence has accumulated rapidly in middle-income settings. Studies in Mexico \citep{hevia2022estimation, alasino2024learning}, South Africa \citep{ardington2021covid}, India \citep{guariso2023impact, singh2024covid}, and Brazil \citep{lichand2022impacts} use in-person assessments of large student samples to document losses of comparable magnitude, while revealing important variation in severity tied to the duration of closure and the quality of remote instruction. Ukraine, an upper-middle-income country, is absent from this literature despite operating a nationally representative primary-school monitoring programme with assessments administered under standardized conditions.

I use two rounds of Ukraine's State Monitoring of Quality of Primary Education
\citep{uceqa2025} to estimate pandemic-related learning losses in mathematics and reading. The programme assessed fourth-grade students in May 2018---before any school disruption---and again in May 2021, after two academic years marked by extended closures and shifts between remote and hybrid instruction. I compare cross-cohort achievement distributions across waves after controlling for region, school type, and student socioeconomic background.

I find that average mathematics scores fell by 0.09 SD between the two cohorts, a decline equivalent to roughly one-fifth of a typical school year \citep{hill2008empirical}. Reading scores fell by a similar magnitude, 0.09--0.10 SD. These losses are smaller than the cross-country averages reported in \citet{betthauser2023systematic} and \citet{patrinos2022analysis}, and I document the mechanisms behind this gap: nearly half of students experienced only one to two months of distance learning, while 13 percent were exposed for three months or longer. I use teacher questionnaires linked to student records to characterize heterogeneity in instructional practices during remote periods.

This paper makes three contributions. First, it provides the first nationally representative estimates of COVID-19 learning losses in Ukrainian primary education. The State Monitoring programme administered identical test instruments under standardized field conditions in both waves and anchored scores to a common scale, which allows direct cross-cohort comparison without the equating assumptions that complicate many international comparisons. Ukraine's experience---an upper-middle-income country with moderately long but variable closure durations---extends the middle-income evidence base to a regional setting that has received little systematic attention. Second, because the design compares cross-sections rather than following the same students, I document sample comparability across waves and assess sensitivity to differential test participation and cohort compositional change. Third, the monitoring programme links student scores to classroom-level teacher questionnaires; I use these data to document heterogeneity in distance-learning exposure and instructional practices, connecting variation in remote-instruction quality to variation in measured losses.

The paper proceeds as follows. Section~\ref{sec:background} describes Ukraine's school closure policies and the remote-learning environment. Section~\ref{sec:data} introduces the monitoring programme, the sample, and the main variables. Section~\ref{sec:results} presents the cross-cohort achievement comparisons. Section~\ref{sec:robustness} reports sensitivity analyses addressing participation and compositional change. Section~\ref{sec:heterogeneity} examines heterogeneity by distance-learning exposure and socioeconomic background. Section~\ref{sec:conclusion} concludes with policy implications.


When governments closed schools to contain COVID-19, they traded an immediate health risk for an educational one. More than 1.6 billion learners across over 190 countries experienced school interruptions during the pandemic \citep{unesco2021year}. The central empirical question is how much those interruptions cost students in learning. The answer is large enough to matter for a generation of workers and for the public investments required to recover.

The evidence from systematic reviews is clear: pandemic-related closures reduced student achievement by economically meaningful amounts, and the losses fell disproportionately on disadvantaged children. \citet{betthauser2023systematic} synthesize 42 studies across 15 countries and document a mean learning deficit of 0.14 standard deviations (SD)---equivalent to roughly one-third of a typical school year---with consistently larger losses among socio-economically disadvantaged students. \citet{patrinos2022analysis} review 36 rigorous evaluations and report a mean shortfall of 0.17 SD, approximately half a grade level, alongside systematic widening of socio-economic gaps. \citet{delacruz2024learning} extend the evidence to 56 studies drawn primarily from developing countries and estimate that each full year of school closure cost approximately 1.1 years of learning; they also identify in-person and phone-based tutoring as effective remediation tools. Together these reviews establish that COVID-19 generated large, unequal, and only partially recovered losses worldwide.

Country-level evidence has accumulated rapidly in middle-income settings. Studies in Mexico \citep{hevia2022estimation, alasino2024learning}, South Africa \citep{ardington2021covid}, India \citep{guariso2023impact, singh2024covid}, and Brazil \citep{lichand2022impacts} use in-person assessments of large student samples to document losses of comparable magnitude, while revealing important variation in severity tied to the duration of closure and the quality of remote instruction. Ukraine, an upper-middle-income country, is absent from this literature despite operating a nationally representative primary-school monitoring programme with assessments administered under standardized conditions.

I use two rounds of Ukraine's State Monitoring of Quality of Primary Education
\citep{uceqa2025} to estimate pandemic-related learning losses in mathematics and reading. The programme assessed fourth-grade students in May 2018---before any school disruption---and again in May 2021, after two academic years marked by extended closures and shifts between remote and hybrid instruction. I compare cross-cohort achievement distributions across waves after controlling for region, school type, and student socioeconomic background.

I find that average mathematics scores fell by 0.09 SD between the two cohorts, a decline equivalent to roughly one-fifth of a typical school year \citep{hill2008empirical}. Reading scores fell by a similar magnitude, 0.09--0.10 SD. These losses are smaller than the cross-country averages reported in \citet{betthauser2023systematic} and \citet{patrinos2022analysis}, and I document the mechanisms behind this gap: nearly half of students experienced only one to two months of distance learning, while 13 percent were exposed for three months or longer. I use teacher questionnaires linked to student records to characterize heterogeneity in instructional practices during remote periods.

This paper makes three contributions. First, it provides the first nationally representative estimates of COVID-19 learning losses in Ukrainian primary education. The State Monitoring programme administered identical test instruments under standardized field conditions in both waves and anchored scores to a common scale, which allows direct cross-cohort comparison without the equating assumptions that complicate many international comparisons. Ukraine's experience---an upper-middle-income country with moderately long but variable closure durations---extends the middle-income evidence base to a regional setting that has received little systematic attention. Second, because the design compares cross-sections rather than following the same students, I document sample comparability across waves and assess sensitivity to differential test participation and cohort compositional change. Third, the monitoring programme links student scores to classroom-level teacher questionnaires; I use these data to document heterogeneity in distance-learning exposure and instructional practices, connecting variation in remote-instruction quality to variation in measured losses.

The paper proceeds as follows. Section~\ref{sec:background} describes Ukraine's school closure policies and the remote-learning environment. Section~\ref{sec:data} introduces the monitoring programme, the sample, and the main variables. Section~\ref{sec:results} presents the cross-cohort achievement comparisons. Section~\ref{sec:robustness} reports sensitivity analyses addressing participation and compositional change. Section~\ref{sec:heterogeneity} examines heterogeneity by distance-learning exposure and socioeconomic background. Section~\ref{sec:conclusion} concludes with policy implications.
